\documentclass[12pt,a4paper]{article}
\usepackage[UTF8]{ctex}
\usepackage{amsmath,amssymb,amsthm}
\usepackage{geometry}
\usepackage{bm}
\usepackage{hyperref}
\usepackage{tikz}
\usetikzlibrary{arrows,arrows.meta,positioning,calc,shapes.geometric}

\geometry{margin=2.5cm}

\title{12.1.7.2 形状闭合的线性规划测试详解}
\author{第12章抓取与操作补充说明}
\date{}

\newtheorem{theorem}{定理}
\newtheorem{definition}{定义}
\newtheorem{example}{例子}
\newtheorem{proposition}{命题}

\begin{document}

\maketitle

\section{章节概述}

12.1.7.2章节介绍如何使用\textbf{线性规划}（Linear Programming, LP）来测试一阶形状闭合。这是一个实用的方法，可以用计算机算法来判断抓取是否产生形状闭合。

\section{形状闭合的数学条件回顾}

\subsection{形状闭合的定义}

\textbf{形状闭合}：如果一组静止约束阻止物体的所有运动，则实现形状闭合。

\textbf{数学表达}：可行twist集合只有原点，即：
\begin{equation}
\mathcal{V} = \{V | F_i^T V \geq 0 \text{ 对所有 } i\} = \{0\}
\end{equation}

\subsection{正张成条件}

形状闭合的充要条件是：接触wrench $F_i$正张成整个wrench空间：
\begin{equation}
\text{pos}(\{F_1, F_2, \ldots, F_j\}) = \mathbb{R}^n
\end{equation}

其中：
\begin{itemize}
\item 对于平面物体：$n = 3$
\item 对于3D物体：$n = 6$
\end{itemize}

\section{线性规划测试的建立}

\subsection{矩阵F的构造}

将$j$个接触wrench排列成矩阵：
\begin{equation}
F = [F_1 \, F_2 \, \cdots \, F_j] \in \mathbb{R}^{n \times j}
\end{equation}

\textbf{注意}：
\begin{itemize}
\item 对于平面物体：$n = 3$，$F_i = [m_{iz} \, f_{ix} \, f_{iy}]^T$
\item 对于3D物体：$n = 6$，$F_i = [m_{ix} \, m_{iy} \, m_{iz} \, f_{ix} \, f_{iy} \, f_{iz}]^T$
\end{itemize}

\subsection{形状闭合的等价条件}

\textbf{条件1}：对于任何外部wrench $F_{\text{ext}} \in \mathbb{R}^n$，存在权重向量$k \in \mathbb{R}^j$，$k \geq 0$，使得：
\begin{equation}
F k + F_{\text{ext}} = 0
\label{eq:form_closure_condition}
\end{equation}

\textbf{关键约束}：$k \geq 0$（所有分量非负）

\textbf{为什么$k$不能是负数？}

\begin{enumerate}
\item \textbf{物理意义}：
   \begin{itemize}
   \item $k_i$表示接触$i$处的\textbf{法向力大小}
   \item 法向力必须是"推"（pushing），不能是"拉"（pulling）
   \item 如果$k_i < 0$，意味着"拉"物体，这在刚体接触中是不可能的
   \item 接触只能提供\textbf{压力}，不能提供\textbf{拉力}
   \end{itemize}

\item \textbf{数学意义}：
   \begin{itemize}
   \item $F k$是接触wrench的\textbf{正线性组合}
   \item 正线性组合要求系数非负（$k_i \geq 0$）
   \item 如果允许$k_i < 0$，就不是正线性组合了
   \end{itemize}

\item \textbf{几何意义}：
   \begin{itemize}
   \item 正线性组合形成\textbf{凸锥}
   \item 如果允许负系数，就是一般的线性组合（整个空间）
   \item 这样就失去了"正张成"的意义
   \end{itemize}
\end{enumerate}

\textbf{物理解释}：
\begin{itemize}
\item $F k$是接触可以提供的wrench（通过压力）
\item $F_{\text{ext}}$是作用在物体上的外部wrench
\item 形状闭合要求：接触可以平衡任何外部wrench
\item 但接触只能提供\textbf{压力}，不能提供拉力
\end{itemize}

\subsection{关键观察}

\textbf{如果$F_{\text{ext}} = 0$}（无外部干扰），条件(\ref{eq:form_closure_condition})变为：
\begin{equation}
F k = 0
\end{equation}

\textbf{关键}：如果存在\textbf{严格正}的系数$k > 0$使得$F k = 0$，则形状闭合成立。

\textbf{为什么需要$k > 0$？}
\begin{itemize}
\item $k > 0$意味着所有接触都"参与"（没有冗余接触）
\item 如果某些$k_i = 0$，说明对应的接触是冗余的
\item 严格正系数确保所有接触都起作用
\end{itemize}

\section{秩条件}

\subsection{必要条件：满秩}

\textbf{命题}：如果$\text{rank}(F) < n$，则物体\textbf{不处于}形状闭合。

\textbf{证明思路}：
\begin{itemize}
\item 如果$\text{rank}(F) < n$，则$F$的列向量不能张成整个$\mathbb{R}^n$
\item 因此不能正张成$\mathbb{R}^n$
\item 形状闭合不成立
\end{itemize}

\textbf{结论}：形状闭合的\textbf{必要条件}是$F$满秩。

\subsection{满秩的检查}

\textbf{方法}：计算矩阵$F$的秩
\begin{equation}
\text{rank}(F) = n \quad \text{（对于形状闭合是必要的）}
\end{equation}

\textbf{注意}：满秩是\textbf{必要但不充分}的条件。

\section{线性规划问题(12.14)}

\subsection{问题表述}

如果$F$是满秩的，形状闭合测试可以表述为以下线性规划问题：

\begin{align}
\text{寻找} \quad & k \\
\text{最小化} \quad & \mathbf{1}^T k \\
\text{使得} \quad & F k = 0 \\
& k_i \geq 1, \quad i = 1, \ldots, j
\label{eq:LP_form_closure}
\end{align}

其中$\mathbf{1}$是$j$维全1向量。

\subsection{约束条件的解释}

\textbf{约束1}：$F k = 0$
\begin{itemize}
\item 这要求接触wrench的线性组合为零
\item 等价于：接触可以产生零wrench（通过内力平衡）
\end{itemize}

\textbf{约束2}：$k_i \geq 1$（对所有$i$）
\begin{itemize}
\item 这确保所有系数\textbf{严格为正}（$k_i > 0$）
\item 因为如果$k_i \geq 1$，则$k_i > 0$
\item 这排除了冗余接触的情况
\end{itemize}

\textbf{目标函数}：$\mathbf{1}^T k = \sum_{i=1}^j k_i$
\begin{itemize}
\item 对于回答"是否形状闭合"这个二元问题，目标函数不是必需的
\item 但包含它可以确保问题\textbf{适定}（well-posed）
\item 某些LP求解器需要目标函数
\end{itemize}

\subsection{测试结果}

\textbf{如果存在解$k$}：
\begin{itemize}
\item 物体处于一阶形状闭合
\item 所有$k_i \geq 1 > 0$，说明所有接触都起作用
\item $F k = 0$，说明接触wrench可以正张成整个空间
\end{itemize}

\textbf{如果不存在解}：
\begin{itemize}
\item 物体不处于形状闭合
\item 可能的原因：
  \begin{itemize}
  \item $F$不是满秩
  \item 即使$F$满秩，也不能正张成整个空间
  \end{itemize}
\end{itemize}

\section{例题12.7的详细分析}

\subsection{问题设置}

\textbf{设置}：
\begin{itemize}
\item 平面物体，中心有孔
\item 两个手指，每个手指接触孔的两条边
\item 总共4个接触点
\item 4个接触wrench：$F_1, F_2, F_3, F_4$
\end{itemize}

\subsection{矩阵F}

给定的矩阵：
\begin{equation}
F = \begin{bmatrix}
0 & 0 & -1 & 2 \\
-1 & 0 & 1 & 0 \\
0 & -1 & 0 & 1
\end{bmatrix}
\end{equation}

\textbf{分析}：
\begin{itemize}
\item $n = 3$（平面问题）
\item $j = 4$（4个接触）
\item $F \in \mathbb{R}^{3 \times 4}$
\end{itemize}

\subsection{秩的检查}

\textbf{检查}：$\text{rank}(F) = 3$

\textbf{验证}：
\begin{itemize}
\item 三个行向量线性无关
\item 或者：三个列向量的某个子集线性无关
\item 因此$F$是满秩的
\end{itemize}

\textbf{结论}：满足形状闭合的必要条件。

\subsection{线性规划求解}

\textbf{问题}：
\begin{align}
\text{最小化} \quad & k_1 + k_2 + k_3 + k_4 \\
\text{使得} \quad & F k = 0 \\
& k_i \geq 1, \quad i = 1, 2, 3, 4
\end{align}

\textbf{解}：$k_1 = 2$，$k_2 = 1$，$k_3 = 2$，$k_4 = 1$

\textbf{验证}：
\begin{align}
F k &= \begin{bmatrix} 0 & 0 & -1 & 2 \\ -1 & 0 & 1 & 0 \\ 0 & -1 & 0 & 1 \end{bmatrix} \begin{bmatrix} 2 \\ 1 \\ 2 \\ 1 \end{bmatrix} \\
&= \begin{bmatrix} 0 \cdot 2 + 0 \cdot 1 + (-1) \cdot 2 + 2 \cdot 1 \\ (-1) \cdot 2 + 0 \cdot 1 + 1 \cdot 2 + 0 \cdot 1 \\ 0 \cdot 2 + (-1) \cdot 1 + 0 \cdot 2 + 1 \cdot 1 \end{bmatrix} \\
&= \begin{bmatrix} -2 + 2 \\ -2 + 2 \\ -1 + 1 \end{bmatrix} = \begin{bmatrix} 0 \\ 0 \\ 0 \end{bmatrix} \quad \checkmark
\end{align}

\textbf{结论}：物体处于形状闭合！

\subsection{物理解释}

\textbf{权重$k$的含义}：
\begin{itemize}
\item $k_1 = 2$：接触1的"贡献"较大
\item $k_2 = 1$：接触2的"贡献"较小
\item $k_3 = 2$：接触3的"贡献"较大
\item $k_4 = 1$：接触4的"贡献"较小
\end{itemize}

\textbf{关键}：所有$k_i > 0$，说明所有接触都起作用，没有冗余。

\section{修改后的例子：不产生形状闭合}

\subsection{新的F矩阵}

如果右手指移动到孔的右下角，新的$F$矩阵为：
\begin{equation}
F = \begin{bmatrix}
0 & 0 & 0 & -2 \\
-1 & 0 & 1 & 0 \\
0 & -1 & 0 & -1
\end{bmatrix}
\end{equation}

\subsection{分析}

\textbf{秩检查}：$\text{rank}(F) = 3$（仍然是满秩）

\textbf{线性规划}：
\begin{itemize}
\item 问题仍然有定义
\item 但\textbf{不存在解}满足所有约束
\end{itemize}

\textbf{结论}：物体\textbf{不处于}形状闭合。

\subsection{为什么没有形状闭合？}

\textbf{原因}：物体可以沿页面\textbf{向下滑动}。

\textbf{分析}：
\begin{itemize}
\item 虽然$F$是满秩的，但不能正张成整个$\mathbb{R}^3$
\item 存在某些方向无法被接触约束阻止
\item 具体来说：向下滑动的方向
\end{itemize}

\section{MATLAB实现}

\subsection{代码说明}

原书提供的MATLAB代码：

\begin{verbatim}
f = [1,1,1,1];  % 目标函数系数
A = [[-1,0,0,0]; [0,-1,0,0]; [0,0,-1,0]; [0,0,0,-1]];  % 不等式约束
b = [-1,-1,-1,-1];  % 不等式约束右端
F = [[0,0,-1,2]; [-1,0,1,0]; [0,-1,0,1]];  % F矩阵
Aeq = F;  % 等式约束
beq = [0,0,0];  % 等式约束右端
k = linprog(f,A,b,Aeq,beq);
\end{verbatim}

\subsection{约束的转换}

\textbf{原始约束}：$k_i \geq 1$

\textbf{标准LP形式}：$A k \leq b$

\textbf{转换}：
\begin{itemize}
\item $k_i \geq 1$等价于$-k_i \leq -1$
\item 因此：$A = -\mathbf{I}$（负单位矩阵）
\item $b = -\mathbf{1}$（全-1向量）
\end{itemize}

\subsection{等式约束}

\textbf{原始约束}：$F k = 0$

\textbf{标准LP形式}：$A_{\text{eq}} k = b_{\text{eq}}$

\textbf{转换}：
\begin{itemize}
\item $A_{\text{eq}} = F$
\item $b_{\text{eq}} = \mathbf{0}$（零向量）
\end{itemize}

\section{为什么$k$不能是负数？详细解释}

\subsection{物理原因}

\textbf{关键理解}：$k_i$表示接触$i$处的\textbf{法向力大小}。

\begin{itemize}
\item \textbf{法向力必须是压力}：
   \begin{itemize}
   \item 接触只能"推"物体，不能"拉"物体
   \item 如果$k_i < 0$，意味着"拉"物体
   \item 这在刚体接触中是不可能的
   \end{itemize}

\item \textbf{接触力的方向}：
   \begin{itemize}
   \item 接触法线$\hat{n}_i$指向物体内部
   \item 接触力必须沿$\hat{n}_i$方向（推入物体）
   \item 不能沿$-\hat{n}_i$方向（拉出物体）
   \end{itemize}
\end{itemize}

\subsection{数学原因}

\textbf{正线性组合的要求}：

形状闭合要求接触wrench \textbf{正张成}整个wrench空间：
\begin{equation}
\text{pos}(\{F_1, \ldots, F_j\}) = \mathbb{R}^n
\end{equation}

\textbf{正张成的定义}：
\begin{equation}
\text{pos}(\{F_1, \ldots, F_j\}) = \left\{\sum_{i=1}^j k_i F_i \middle| k_i \geq 0\right\}
\end{equation}

\textbf{关键}：系数$k_i$必须\textbf{非负}！

如果允许$k_i < 0$：
\begin{itemize}
\item 就不是正线性组合了
\item 变成一般的线性组合
\item 失去了"正张成"的意义
\end{itemize}

\subsection{几何原因}

\begin{itemize}
\item \textbf{正线性组合}：形成以原点为顶点的\textbf{凸锥}
\item \textbf{一般线性组合}：形成整个空间（如果$F_i$线性无关）
\item 形状闭合需要的是\textbf{凸锥}，不是整个空间
\end{itemize}

\subsection{具体例子}

\textbf{例子}：两个接触wrench $F_1$和$F_2$

\begin{itemize}
\item \textbf{如果$k_1, k_2 \geq 0$}：
   \begin{itemize}
   \item $k_1 F_1 + k_2 F_2$形成凸锥
   \item 只能表示某些方向
   \end{itemize}

\item \textbf{如果允许$k_1, k_2 < 0$}：
   \begin{itemize}
   \item $k_1 F_1 + k_2 F_2$可以表示任何方向
   \item 但这不是"正张成"
   \end{itemize}
\end{itemize}

\section{为什么需要$k_i \geq 1$而不是$k_i > 0$？}

\subsection{数学原因}

\textbf{问题}：为什么用$k_i \geq 1$而不是$k_i > 0$？

\textbf{答案}：
\begin{itemize}
\item 线性规划通常只支持\textbf{弱不等式}（$\geq$或$\leq$），不支持严格不等式（$>$或$<$）
\item $k_i \geq 1$确保$k_i > 0$（因为$1 > 0$）
\item 这是处理严格正约束的标准方法
\end{itemize}

\subsection{等价性证明}

\textbf{关键观察}：
\begin{itemize}
\item 如果存在$k > 0$（严格正）使得$F k = 0$
\item 那么存在$\tilde{k} = \alpha k$（$\alpha > 0$足够大）使得$\tilde{k}_i \geq 1$
\item 因为如果$k_i > 0$，我们可以选择$\alpha \geq \frac{1}{\min_i k_i}$，使得$\tilde{k}_i = \alpha k_i \geq 1$
\end{itemize}

\textbf{因此}：$k_i \geq 1$和$k_i > 0$在这个问题中是等价的。

\section{算法步骤总结}

\subsection{测试流程}

\begin{enumerate}
\item \textbf{步骤1}：构造矩阵$F = [F_1 \, F_2 \, \cdots \, F_j]$

\item \textbf{步骤2}：检查秩
   \begin{itemize}
   \item 如果$\text{rank}(F) < n$：\textbf{不是}形状闭合，结束
   \item 如果$\text{rank}(F) = n$：继续
   \end{itemize}

\item \textbf{步骤3}：求解线性规划问题(12.14)
   \begin{itemize}
   \item 如果存在解$k$：\textbf{是}形状闭合
   \item 如果不存在解：\textbf{不是}形状闭合
   \end{itemize}
\end{enumerate}

\section{关键理解点}

\subsection{为什么用线性规划？}

\begin{enumerate}
\item \textbf{可计算性}：
   \begin{itemize}
   \item 线性规划有高效的求解算法
   \item 可以在多项式时间内求解
   \end{itemize}

\item \textbf{实用性}：
   \begin{itemize}
   \item 可以用标准软件（如MATLAB）求解
   \item 适合自动化测试
   \end{itemize}

\item \textbf{精确性}：
   \begin{itemize}
   \item 给出明确的"是"或"否"答案
   \item 不依赖图形方法
   \end{itemize}
\end{enumerate}

\subsection{与正张成条件的关系}

\textbf{正张成条件}：$\text{pos}(\{F_1, \ldots, F_j\}) = \mathbb{R}^n$

\textbf{线性规划条件}：存在$k > 0$使得$F k = 0$

\textbf{等价性}：
\begin{itemize}
\item 如果$F k = 0$且$k > 0$，则$F_i$可以正张成$\mathbb{R}^n$
\item 反过来也成立
\item 因此两个条件是等价的
\end{itemize}

\section{实际应用}

\subsection{抓取验证}

\begin{itemize}
\item 设计抓取后，可以用LP测试验证是否产生形状闭合
\item 快速、准确、可自动化
\end{itemize}

\subsection{抓取优化}

\begin{itemize}
\item 可以结合优化方法
\item 寻找产生形状闭合的接触位置
\item 使用LP测试作为约束条件
\end{itemize}

\section{总结}

\subsection{关键公式}

\textbf{线性规划问题(12.14)}：
\begin{align}
\text{最小化} \quad & \mathbf{1}^T k \\
\text{使得} \quad & F k = 0 \\
& k_i \geq 1, \quad i = 1, \ldots, j
\end{align}

\subsection{测试结果}

\begin{itemize}
\item \textbf{如果存在解}：形状闭合
\item \textbf{如果不存在解}：不是形状闭合
\end{itemize}

\subsection{前提条件}

\begin{itemize}
\item $F$必须是满秩的（$\text{rank}(F) = n$）
\item 否则直接判断为不是形状闭合
\end{itemize}

\subsection{优势}

\begin{enumerate}
\item \textbf{可计算}：可以用标准算法求解
\item \textbf{精确}：给出明确答案
\item \textbf{实用}：适合自动化测试
\item \textbf{通用}：适用于任何数量的接触
\end{enumerate}

\end{document}

